\documentclass[11pt,a4paper]{report}
\usepackage[utf8]{inputenc}
\usepackage[T1]{fontenc}
\usepackage[italian]{babel}
\usepackage{lmodern}
\usepackage[margin=1.8cm]{geometry}
\usepackage{amsmath,amssymb,amsthm,mathtools}
\usepackage{booktabs}
\usepackage{tabularx}
\usepackage{xcolor}
\usepackage{enumitem}
\usepackage[most]{tcolorbox}
\usepackage{hyperref}

\hypersetup{
    colorlinks=true,
    linkcolor=blue!70!black,
    citecolor=blue!70!black,
    urlcolor=blue!70!black
}

% Palette Colori
\definecolor{primaryblue}{RGB}{24, 76, 120}
\definecolor{boxbg}{RGB}{248, 250, 252}
\definecolor{boxframe}{RGB}{52, 116, 170}
\definecolor{intbg}{RGB}{240, 253, 244}
\definecolor{intframe}{RGB}{34, 139, 34}
\definecolor{usebg}{RGB}{238, 246, 255}
\definecolor{useframe}{RGB}{41, 128, 185}
\definecolor{demobg}{RGB}{250, 245, 255}
\definecolor{demoframe}{RGB}{128, 0, 128}
\definecolor{oralbg}{RGB}{255, 251, 235}
\definecolor{oralframe}{RGB}{217, 119, 6}
\definecolor{tipbg}{RGB}{254, 242, 242}
\definecolor{tipframe}{RGB}{185, 28, 28}

% Box Teorema Principale
\newtcolorbox{teoremabox}[2][]{
    enhanced, breakable,
    colback=boxbg, colframe=boxframe,
    fonttitle=\bfseries\color{white},
    title={#2}, arc=2mm, boxrule=1.2pt,
    left=3.5mm, right=3.5mm, top=3mm, bottom=3mm,
    before skip=10pt, after skip=6pt, #1
}

% Sotto-Box 1: Cosa fa (Intuizione)
\newtcolorbox{cosafa}[1][Cosa Fa (L'Intuizione in parole semplici)]{
    enhanced, breakable,
    colback=intbg, colframe=intframe,
    fonttitle=\bfseries\small\color{white},
    title={$\blacktriangleright$ #1}, arc=1.5mm, boxrule=0.8pt,
    left=2.5mm, right=2.5mm, top=2mm, bottom=2mm,
    before skip=4pt, after skip=4pt
}

% Sotto-Box 2: Dimostrazione
\newtcolorbox{dimostrazione}[1][Dimostrazione Passo-Passo (Snella \& da Lavagna)]{
    enhanced, breakable,
    colback=demobg, colframe=demoframe,
    fonttitle=\bfseries\small\color{white},
    title={$\blacktriangleright$ #1}, arc=1.5mm, boxrule=0.8pt,
    left=2.5mm, right=2.5mm, top=2mm, bottom=2mm,
    before skip=4pt, after skip=4pt
}

% Sotto-Box 3: Come dirla al prof
\newtcolorbox{oralbox}[1][Come Dirla al Prof (Discorso da 30L)]{
    enhanced, breakable,
    colback=oralbg, colframe=oralframe,
    fonttitle=\bfseries\small\color{white},
    title={$\blacktriangleright$ #1}, arc=1.5mm, boxrule=0.8pt,
    left=2.5mm, right=2.5mm, top=2mm, bottom=2mm,
    before skip=4pt, after skip=4pt
}

% Sotto-Box 4: Trucco per ricordare
\newtcolorbox{tipbox}[1][Suggerimento per Ricordarla Facilmente]{
    enhanced, breakable,
    colback=tipbg, colframe=tipframe,
    fonttitle=\bfseries\small\color{white},
    title={$\blacktriangleright$ #1}, arc=1.5mm, boxrule=0.8pt,
    left=2.5mm, right=2.5mm, top=2mm, bottom=2mm,
    before skip=4pt, after skip=4pt
}

\begin{document}

\title{\textbf{\LARGE I 16 Teoremi ``Puri'' d'Esame}\\[0.3cm]
\large Enunciati Letterali, Glossario, Dimostrazioni Snelle e Guida all'Esposizione Orale}
\author{\textbf{Francesco Castaldi} (\texttt{francesco.castaldi4@studio.unibo.it})\\
Corso di Laurea in Informatica --- Università di Bologna\\
Insegnamento: \emph{Modelli Probabilistici (518512)} --- Docente: Prof. Salvatore Federico}
\date{Anno Accademico 2025/2026 --- Appello di Settembre 2026}

\maketitle
\tableofcontents
\newpage

% ============================================================
\chapter*{Glossario: Parole Chiave Spiegate in Modo Semplice}
\addcontentsline{toc}{chapter}{Glossario: Parole Chiave Spiegate in Modo Semplice}
% ============================================================

\begin{itemize}[leftmargin=*, itemsep=6pt]
    \item \textbf{Spazio di Probabilità $(\Omega, \mathcal{F}, \mathbb{P})$}: L'universo dell'esperimento casuale. $\Omega$ è l'insieme di tutti i possibili esiti; $\mathcal{F}$ è la collezione degli eventi osservabili (le domande a cui possiamo rispondere con certezza); $\mathbb{P}$ assegna a ciascun evento un valore tra 0 e 1.
    \item \textbf{Variabile Aleatoria (v.a.) $X$}: Una funzione deterministica $X: \Omega \to E$ che converte ogni esito dell'universo $\Omega$ in un valore numerico (o stato) misurabile.
    \item \textbf{i.i.d. (Indipendenti e Identicamente Distribuite)}: Variabili generate da esperimenti ripetuti nelle identiche condizioni, in cui il passato non altera il futuro e la legge di probabilità è sempre la medesima.
    \item \textbf{Valore Atteso $\mathbb{E}[X]$}: La media ponderata dei valori assunti da $X$, pesati con le rispettive probabilità.
    \item \textbf{Convergenza Monotona $X_n \uparrow X$}: Una successione di variabili che cresce punto per punto ($X_1(\omega) \le X_2(\omega) \le \dots$) convergendo al tetto $X(\omega)$.
    \item \textbf{Quasi Certamente (q.c. / a.s.)}: Un evento che accade con probabilità $1$.
    \item \textbf{Misurabilità rispetto a una $\sigma$-algebra $\mathcal{G}$}: Significa che il valore della variabile è interamente deducibile dall'informazione racchiusa in $\mathcal{G}$.
    \item \textbf{Tempo di Arresto $\tau$}: Una variabile aleatoria a valori interi $\ge 0$ tale che l'evento $\{\tau \le n\}$ dipende solo dai dati osservati fino al tempo $n$ (senza conoscere il futuro).
    \item \textbf{Catena di Markov}: Un processo stocastico senza memoria storica: la distribuzione del passo futuro $X_{n+1}$ dipende unicamente dallo stato corrente $X_n$.
    \item \textbf{Stato Ricorrente vs Transiente}: Uno stato $x$ è \emph{ricorrente} se partendo da $x$ la probabilità di ritornarvi prima o poi è pari a $1$; è \emph{transiente} se tale probabilità è strettamente minore di $1$.
    \item \textbf{Distribuzione Stazionaria $\pi$}: Una misura di probabilità invariante per la matrice di transizione ($\pi P = \pi$), che descrive l'equilibrio statistico a lungo termine.
\end{itemize}

\newpage

% ============================================================
\chapter{Convergenze e Leggi dei Grandi Numeri (Cap. 2)}
% ============================================================

\begin{teoremabox}{Teorema 2.4.1 --- Teorema di Estensione di Kolmogorov}
Data una famiglia di leggi di probabilità $\{\mu_i\}_{i=1}^N$ su spazi discreti $\{E_i\}_{i=1}^N$ (con $N \le \infty$), è sempre possibile costruire uno spazio di probabilità $(\Omega, \mathcal{F}, \mathbb{P})$ e una successione di v.a. indipendenti $\{X_i\}_{i=1}^N$ tali che $\mathcal{L}(X_i) = \mu_i$.

\begin{cosafa}
Garantisce la fondatezza matematica dei modelli probabilistici: assicura che possiamo sempre definire infinite variabili aleatorie indipendenti senza generare contraddizioni negli assiomi.
\end{cosafa}

\begin{dimostrazione}
Si definisce lo spazio prodotto $\Omega = \prod_{i=1}^\infty E_i$. Sugli insiemi cilindrici finiti $C = \{ \omega : \omega_1 = x_1, \dots, \omega_k = x_k \}$, si pone la misura prodotto $\mathbb{P}(C) = \prod_{i=1}^k \mu_i(\{x_i\})$. Il Teorema di Carathéodory garantisce l'estensione unica di $\mathbb{P}$ all'intera $\sigma$-algebra prodotto $\mathcal{F} = \bigotimes E_i$.
\end{dimostrazione}

\begin{oralbox}
``Professore, questo risultato assicura l'esistenza rigorosa di sequenze infinite di lanci o prove indipendenti partendo semplicemente dalle singole distribuzioni marginali.''
\end{oralbox}

\begin{tipbox}
Ricorda la parola chiave: \textbf{Spazio Prodotto + Misura Prodotto sui Cilindri}.
\end{tipbox}
\end{teoremabox}

\vspace{0.3cm}

\begin{teoremabox}{Teorema 2.7.6 --- Teorema di Convergenza Monotona (Beppo-Levi)}
Sia $(X_n)_{n \ge 0}$ una successione di variabili aleatorie discrete tali che $X_n \ge 0$ e $X_n \uparrow X$ q.c. Allora:
\[
\lim_{n \to \infty} \mathbb{E}[X_n] = \mathbb{E}[X].
\]

\begin{cosafa}
Permette di scambiare l'operatore di limite con il valore atteso quando le variabili sono positive e crescono monotonicamente ad ogni passo.
\end{cosafa}

\begin{dimostrazione}
\begin{enumerate}
    \item Poiché $X_n \le X_{n+1} \le X$, per monotonia del valore atteso la successione numerica $\mathbb{E}[X_n]$ è non decrescente e limitata superiormente da $\mathbb{E}[X]$, dunque ammette limite $L \le \mathbb{E}[X]$.
    \item Scrivendo il valore atteso come somma di serie su spazi discreti $\mathbb{E}[X_n] = \sum_x x \mathbb{P}(X_n = x)$, lo scambio tra serie a termini non negativi crescenti e limite è garantito dal teorema classico di Beppo-Levi sulle serie reali. Dunque $L = \mathbb{E}[X]$.
\end{enumerate}
\end{dimostrazione}

\begin{oralbox}
``Trattandosi di spazi discreti, il valore atteso è una somma di serie; la positività e la monotonicità assicurano che il limite passi direttamente sotto il segno di sommatoria.''
\end{oralbox}

\begin{tipbox}
Pensa all'analogia di un salvadanaio: se metti solo monete positive e il capitale sale sempre, la media finale coincide con il limite delle medie parziali.
\end{tipbox}
\end{teoremabox}

\newpage

\begin{teoremabox}{Teorema 2.7.10 --- Teorema di Convergenza Dominata di Lebesgue}
Sia $(X_n)_{n \ge 0}$ una successione di v.a. discrete con $X_n \to X$ q.c. Se esiste $G \in L^1$ ($\mathbb{E}[G] < \infty$) tale che $|X_n| \le G$ per ogni $n$, allora $X \in L^1$ e:
\[
\lim_{n \to \infty} \mathbb{E}[X_n] = \mathbb{E}[X].
\]

\begin{cosafa}
Fornisce la condizione universale per scambiare limite e media: non serve la monotonicità, basta che la sequenza non esploda e sia ``schiacciata'' sotto un tetto integrabile $G$.
\end{cosafa}

\begin{dimostrazione}
Poiché $|X_n| \le G$, definiamo due successioni non negative: $Y_n = G + X_n \ge 0$ e $Z_n = G - X_n \ge 0$. Applichiamo il Lemma di Fatou a entrambe:
\begin{enumerate}
    \item $\mathbb{E}[G + X] \le \liminf \mathbb{E}[G + X_n] \implies \mathbb{E}[X] \le \liminf \mathbb{E}[X_n]$.
    \item $\mathbb{E}[G - X] \le \liminf \mathbb{E}[G - X_n] \implies \mathbb{E}[X] \ge \limsup \mathbb{E}[X_n]$.
\end{enumerate}
Unendo le disuguaglianze: $\limsup \mathbb{E}[X_n] \le \mathbb{E}[X] \le \liminf \mathbb{E}[X_n]$, da cui $\lim_{n \to \infty} \mathbb{E}[X_n] = \mathbb{E}[X]$.
\end{dimostrazione}

\begin{oralbox}
``Dimostriamo la convergenza dominata applicando Fatou a $G+X_n$ e $G-X_n$, costringendo liminf e limsup a coincidere con $\mathbb{E}[X]$.''
\end{oralbox}

\begin{tipbox}
Formula chiave da ricordare: \textbf{Doppia applicazione di Fatou con $G \pm X_n$}.
\end{tipbox}
\end{teoremabox}

\vspace{0.3cm}

\begin{teoremabox}{Teorema 2.7.11 --- Versioni Condizionate dei Teoremi di Convergenza}
Sia $Y$ una v.a. discreta. Valgono le estensioni al valore atteso condizionato:
\begin{enumerate}
    \item \textbf{Monotona condizionata}: Se $0 \le X_n \uparrow X$ q.c. $\implies \mathbb{E}[X_n \mid Y] \uparrow \mathbb{E}[X \mid Y]$.
    \item \textbf{Fatou condizionato}: Se $X_n \ge 0$ q.c. $\implies \mathbb{E}[\liminf X_n \mid Y] \le \liminf \mathbb{E}[X_n \mid Y]$.
    \item \textbf{Dominata condizionata}: Se $X_n \to X$ q.c. e $|X_n| \le Z \in L^1 \implies \lim \mathbb{E}[X_n \mid Y] = \mathbb{E}[X \mid Y]$.
\end{enumerate}

\begin{cosafa}
Conferma che tutte le proprietà di passaggio al limite rimangono valide anche all'interno delle previsioni condizionali rispetto a dati noti.
\end{cosafa}

\begin{dimostrazione}
Su ogni atomo $\{Y = y\}$ con $\mathbb{P}(Y=y) > 0$, il valore atteso condizionato si riduce a un'ordinaria media pesata sulla misura condizionale $\mathbb{P}(\cdot \mid Y=y)$. Applicando i teoremi classici su ciascun atomo separatamente, la proprietà vale su tutto $\Omega$.
\end{dimostrazione}

\begin{oralbox}
``Le versioni condizionate discendono direttamente dai teoremi standard valutati sulla misura di probabilità condizionata $\mathbb{P}(\cdot \mid Y=y)$ su ciascun atomo.''
\end{oralbox}
\end{teoremabox}

\newpage

\begin{teoremabox}{Teorema 2.8.1 --- Legge Debole dei Grandi Numeri (WLLN)}
Sia $(X_n)_{n \ge 1}$ i.i.d. con $\mathbb{E}[X_1] = m$ e $\mathrm{Var}(X_1) = \sigma^2 < \infty$. Posto $S_n = \sum_{k=1}^n X_k$, per ogni $\epsilon > 0$:
\[
\lim_{n \to \infty} \mathbb{P}\left(\left| \frac{S_n}{n} - m \right| \ge \epsilon\right) = 0.
\]

\begin{cosafa}
Dimostra che la media aritmetica di $n$ campioni indipendenti converge in probabilità alla media teorica $m$.
\end{cosafa}

\begin{dimostrazione}
\begin{enumerate}
    \item Calcoliamo media e varianza della media campionaria:
    \[
    \mathbb{E}\left[\frac{S_n}{n}\right] = m, \quad \mathrm{Var}\left(\frac{S_n}{n}\right) = \frac{1}{n^2} \sum_{k=1}^n \mathrm{Var}(X_k) = \frac{n\sigma^2}{n^2} = \frac{\sigma^2}{n}.
    \]
    \item Applichiamo la Disuguaglianza di Chebyshev:
    \[
    \mathbb{P}\left(\left| \frac{S_n}{n} - m \right| \ge \epsilon\right) \le \frac{\mathrm{Var}(S_n/n)}{\epsilon^2} = \frac{\sigma^2}{n \epsilon^2}.
    \]
    \item Facendo tendere $n \to \infty$, il membro destro tende a $0$.
\end{enumerate}
\end{dimostrazione}

\begin{oralbox}
``Calcoliamo la varianza della media campionaria che scala come $\sigma^2/n$ e applichiamo Chebyshev: al tendere di $n$ all'infinito la probabilità decade a zero.''
\end{oralbox}

\begin{tipbox}
Schema mentale: \textbf{Varianza $\sigma^2/n$ + Chebyshev}.
\end{tipbox}
\end{teoremabox}

\vspace{0.3cm}

\begin{teoremabox}{Teorema 2.8.2 --- Legge Forte di Borel (SLLN per v.a. limitate)}
Sia $(X_n)_{n \ge 1}$ i.i.d. limitata ($|X_n| \le C < \infty$) con $\mathbb{E}[X_1] = m$. Allora:
\[
\frac{S_n}{n} \longrightarrow m \quad \text{quasi certamente per } n \to \infty.
\]

\begin{cosafa}
Rafforza la legge debole: la convergenza della media empirica al valore atteso teorico avviene con certezza assoluta (probabilità $1$) su quasi ogni traiettoria.
\end{cosafa}

\begin{dimostrazione}
Assumiamo senza perdita di generalità $m=0$.
\begin{enumerate}
    \item Calcoliamo il momento quarto $\mathbb{E}[S_n^4] = \mathbb{E}\left[\left(\sum X_i\right)^4\right]$. Sviluppando il prodotto, per indipendenza e media nulla sopravvivono solo i termini $\mathbb{E}[X_i^4] \le C^4$ e $\mathbb{E}[X_i^2 X_j^2] \le C^4$. I termini misti sono $\binom{n}{2}\binom{4}{2} = 3n(n-1) \le 3n^2$. Dunque $\mathbb{E}[S_n^4] \le K n^2$.
    \item Per Chebyshev con momento quarto: $\mathbb{P}(|S_n/n| \ge \epsilon) \le \frac{\mathbb{E}[S_n^4]}{n^4 \epsilon^4} \le \frac{K}{n^2 \epsilon^4}$.
    \item Poiché $\sum_{n=1}^\infty \frac{1}{n^2} < \infty$, per il Primo Lemma di Borel-Cantelli l'evento $\{|S_n/n| \ge \epsilon\}$ si verifica solo un numero finito di volte q.c., implicando la convergenza quasi certa.
\end{enumerate}
\end{dimostrazione}

\begin{oralbox}
``Sfruttiamo il momento quarto $\mathbb{E}[S_n^4] = O(n^2)$: dividendo per $n^4$ la serie delle probabilità converge come $\sum 1/n^2$ e Borel-Cantelli conclude.''
\end{oralbox}

\begin{tipbox}
Ricorda: \textbf{Momento $4^\circ \implies O(n^2)/n^4 = 1/n^2 \implies$ Borel-Cantelli}.
\end{tipbox}
\end{teoremabox}

\newpage

% ============================================================
\chapter{Passeggiate Aleatorie e Borel-Cantelli (Cap. 5)}
% ============================================================

\begin{teoremabox}{Teorema 5.1.2 --- Legge 0-1 di Kolmogorov}
Sia $(X_n)_{n \ge 1}$ una successione di v.a. discrete indipendenti e sia $\mathcal{T} = \bigcap_{n=1}^\infty \sigma(X_n, X_{n+1}, \dots)$ la $\sigma$-algebra coda. Se $A \in \mathcal{T}$, allora:
\[
\mathbb{P}(A) = 0 \quad \text{oppure} \quad \mathbb{P}(A) = 1.
\]

\begin{cosafa}
Stabilisce che gli eventi asintotici di lunghissimo periodo (che non dipendono dai primi $N$ passi finiti) non ammettono incertezza intermedia: sono certi oppure impossibili.
\end{cosafa}

\begin{dimostrazione}
\begin{enumerate}
    \item Sia $A \in \mathcal{T}$. Poiché $\mathcal{T} \subseteq \sigma(X_{n+1}, X_{n+2}, \dots)$, l'evento $A$ è indipendente dalla storia iniziale $\mathcal{F}_n = \sigma(X_1, \dots, X_n)$ per ogni $n$.
    \item Per passaggio al limite, $A$ è indipendente dall'intera algebra $\bigcup \mathcal{F}_n$ e quindi dalla $\sigma$-algebra generata $\sigma(X_1, X_2, \dots)$.
    \item Ma $A \in \mathcal{T} \subseteq \sigma(X_1, X_2, \dots)$, quindi $A$ è indipendente da se stesso:
    \[
    \mathbb{P}(A) = \mathbb{P}(A \cap A) = \mathbb{P}(A) \cdot \mathbb{P}(A) = \mathbb{P}(A)^2.
    \]
    \item L'equazione $p = p^2$ ammette unicamente le soluzioni $p=0$ e $p=1$.
\end{enumerate}
\end{dimostrazione}

\begin{oralbox}
``Un evento coda è indipendente da qualsiasi orizzonte finito $n$, dunque è indipendente dall'intera successione e da se stesso, imponendo $\mathbb{P}(A) = \mathbb{P}(A)^2 \implies \mathbb{P}(A) \in \{0,1\}$.''
\end{oralbox}

\begin{tipbox}
Trucco: \textbf{Indipendenza da se stesso $\implies \mathbb{P}(A) = \mathbb{P}(A)^2$}.
\end{tipbox}
\end{teoremabox}

\vspace{0.3cm}

\begin{teoremabox}{Teorema 5.2.2 --- Limite Superiore Infinito per Random Walk}
Per una passeggiata aleatoria simmetrica su $\mathbb{Z}$ con passi i.i.d. non degeneri:
\[
\sup_{n \ge 0} |X_n| = \infty \quad \text{quasi certamente}.
\]

\begin{cosafa}
Garantisce che una passeggiata aleatoria equa non rimane confinata in una gabbia limitata: continuerà a oscillare esplorando spazi arbitrariamente grandi.
\end{cosafa}

\begin{dimostrazione}
L'evento $A = \{\sup_{n \ge 0} |X_n| < \infty\}$ appartiene alla $\sigma$-algebra coda $\mathcal{T}$, dunque per la Legge 0-1 di Kolmogorov $\mathbb{P}(A) \in \{0,1\}$. Se fosse $\mathbb{P}(A) = 1$, la passeggiata sarebbe uniformemente limitata, contraddicendo la varianza $\mathrm{Var}(X_n) = n\sigma^2 \to \infty$. Dunque necessariamente $\mathbb{P}(A) = 0$, da cui $\mathbb{P}(\sup |X_n| = \infty) = 1$.
\end{dimostrazione}

\begin{oralbox}
``L'evento di essere limitata è un evento coda con probabilità 0 o 1; poiché la varianza cresce come $n$, non può essere limitata, quindi il sup è infinito q.c.''
\end{oralbox}
\end{teoremabox}

\newpage

% ============================================================
\chapter{Catene di Markov su Spazi Discreti (Cap. 7)}
% ============================================================

\begin{teoremabox}{Teorema 7.3.1 --- Esistenza e Unicità in Legge delle Catene di Markov}
Dato uno spazio discreto $E$, una legge iniziale $\mu_0$ su $E$ e un nucleo di transizione stocastico $P = (p_{ij})_{i,j \in E}$, esiste una Catena di Markov $X = (X_n)_{n \ge 0}$ con distribuzione iniziale $\mu_0$ e nucleo $P$, ed è unica in legge sulle traiettorie $E^{\mathbb{N}}$.

\begin{cosafa}
Assicura che specificare il punto di partenza e la matrice delle probabilità di salto ad un passo individua in modo univoco l'intero processo temporale.
\end{cosafa}

\begin{dimostrazione}
Definiamo le distribuzioni finito-dimensionali su cilindri:
\[
\mathbb{P}(X_0 = i_0, X_1 = i_1, \dots, X_n = i_n) = \mu_0(i_0) p_{i_0 i_1} p_{i_1 i_2} \cdots p_{i_{n-1} i_n}.
\]
Tale famiglia soddisfa le condizioni di consistenza di Kolmogorov. Per il Teorema di Estensione di Kolmogorov esiste un'unica misura di probabilità $\mathbb{P}$ su $E^{\mathbb{N}}$.
\end{dimostrazione}

\begin{oralbox}
``Costruiamo le distribuzioni congiunte sui cilindri moltiplicando la misura iniziale per le probabilità di transizione a catena; Kolmogorov garantisce esistenza e unicità globale.''
\end{oralbox}
\end{teoremabox}

\vspace{0.3cm}

\begin{teoremabox}{Teorema 7.4.1 --- Rappresentazione tramite Equazioni alle Differenze Stocastiche}
Sia $E$ finito. Ogni catena di Markov con matrice di transizione razionale può essere espressa come:
\[
X_{n+1} = f(X_n, \xi_{n+1}),
\]
dove $(\xi_n)_{n \ge 1}$ è una successione i.i.d. di variabili uniformi discrete e $f$ è deterministica.

\begin{cosafa}
Dimostra il principio su cui si basano le simulazioni informatiche (es. simulatore Monte Carlo Python nel nostro modello SIR): il passo successivo è una funzione deterministica dello stato attuale più un numero casuale indipendente.
\end{cosafa}

\begin{dimostrazione}
Sia $M$ il minimo comune multiplo dei denominatori delle probabilità $p_{ij} = k_{ij}/M$. Suddividiamo l'insieme $\{1, \dots, M\}$ in intervalli disgiunti $I_{ij}$ di ampiezza $k_{ij}$. Definiamo $f(i, u) = j$ se $u \in I_{ij}$. Se $\xi \sim \mathrm{Unif}(\{1, \dots, M\})$, allora $\mathbb{P}(f(i, \xi) = j) = |I_{ij}|/M = p_{ij}$.
\end{dimostrazione}

\begin{oralbox}
``È l'algoritmo di campionamento inverso: partizioniamo l'intervallo con le probabilità cumulate e mappiamo il rumore uniforme indipendente sul nuovo stato.''
\end{oralbox}

\begin{tipbox}
Collegamento all'esame: \textbf{È il listing Python del simulatore SIR!}
\end{tipbox}
\end{teoremabox}

\newpage

\begin{teoremabox}{Teorema 7.6.3 --- Proprietà Forte di Markov (Strong Markov Property)}
Sia $X$ una catena di Markov omogenea con matrice $P$ e sia $\tau$ un tempo di arresto con $\mathbb{P}(\tau < \infty) > 0$. Allora, condizionatamente a $\{\tau < \infty\}$, il processo post-arresto $Y_n = X_{\tau + n}$ è ancora una catena di Markov omogenea con matrice $P$, indipendente dalla storia passata $\mathcal{F}_\tau$.

\begin{cosafa}
Estende la perdita di memoria: la catena ``riparte da zero'' dimenticando il passato non solo a tempi deterministici prefissati, ma anche ad ogni istante casuale deciso da una regola di stop.
\end{cosafa}

\begin{dimostrazione}
Scomponiamo l'evento condizionante sulla partizione numerabile $\{\tau = k\}$:
\[
\mathbb{P}(X_{\tau+1} = j \mid X_\tau = i, A \cap \{\tau = k\}) = \mathbb{P}(X_{k+1} = j \mid X_k = i, A \cap \{\tau = k\}).
\]
Poiché $A \cap \{\tau = k\} \in \mathcal{F}_k$, per la Proprietà di Markov semplice al tempo fisso $k$ il passato scompare, lasciando $p_{ij}$. Sommando su tutti i $k$, otteniamo l'invarianza anche al tempo d'arresto.
\end{dimostrazione}

\begin{oralbox}
``Disintegriamo il tempo d'arresto sui singoli istanti fissi $\{\tau = k\}$; poiché su ciascuno vale la Markov debole, ricomponendo la somma otteniamo la proprietà forte.''
\end{oralbox}

\begin{tipbox}
Formula chiave: \textbf{Disintegrazione su $\{\tau = k\}$ + Markov semplice su ogni $k$}.
\end{tipbox}
\end{teoremabox}

\vspace{0.3cm}

\begin{teoremabox}{Teorema 7.7.4 --- Probabilità di Ritorno Molteplice $F(x,x)^k$}
Sia $\tau_x^{(k)}$ l'istante del $k$-esimo ritorno allo stato $x$, e sia $F(x,x) = \mathbb{P}_x(\tau_x < \infty)$ la probabilità di primo ritorno. Allora per ogni $k \ge 1$:
\[
\mathbb{P}_x(\tau_x^{(k)} < \infty) = [F(x,x)]^k.
\]

\begin{cosafa}
Dimostra che ogni ciclo di ritorno è identico al primo: se la probabilità di tornare una volta è $F$, tornare $k$ volte è come vincere $k$ estrazioni indipendenti consecutive ($F^k$).
\end{cosafa}

\begin{dimostrazione}
Per induzione su $k$. Per $k=1$ è la definizione. Supponiamo vero per $k-1$. Il $k-1$-esimo ritorno $\tau_x^{(k-1)}$ è un tempo di arresto finito. Per la \textbf{Proprietà Forte di Markov}, il processo riparte dallo stato $x$:
\[
\mathbb{P}_x(\tau_x^{(k)} < \infty) = \mathbb{P}_x(\tau_x^{(k-1)} < \infty) \cdot \mathbb{P}_x(\tau_x < \infty) = F(x,x)^{k-1} \cdot F(x,x) = F(x,x)^k.
\]
\end{dimostrazione}

\begin{oralbox}
``Ogni ritorno è un tempo di arresto: la proprietà forte resetta la catena e moltiplica la probabilità $F(x,x)$ ad ogni nuovo ciclo.''
\end{oralbox}
\end{teoremabox}

\newpage

\begin{teoremabox}{Teorema 7.7.10 --- Solidarietà di Classe della Ricorrenza}
Sia $C$ una classe comunicante ($x \leftrightarrow y$ per ogni $x,y \in C$). Se uno stato $x \in C$ è ricorrente, allora tutti gli stati di $C$ sono ricorrenti.

\begin{cosafa}
Stabilisce che la ricorrenza è una proprietà di classe: in un gruppo di stati reciprocamente accessibili, o sono tutti ricorrenti o sono tutti transienti.
\end{cosafa}

\begin{dimostrazione}
\begin{enumerate}
    \item Sia $x$ ricorrente, dunque $\sum_{n=0}^\infty p_{xx}^{(n)} = \infty$.
    \item Poiché $x \leftrightarrow y$, esistono interi $m, r > 0$ tali che $p_{yx}^{(m)} > 0$ e $p_{xy}^{(r)} > 0$.
    \item Per le equazioni di Chapman-Kolmogorov:
    \[
    p_{yy}^{(m+n+r)} \ge p_{yx}^{(m)} p_{xx}^{(n)} p_{xy}^{(r)}.
    \]
    \item Sommando su tutti gli $n$:
    \[
    \sum_{n=0}^\infty p_{yy}^{(n)} \ge p_{yx}^{(m)} p_{xy}^{(r)} \sum_{n=0}^\infty p_{xx}^{(n)} = p_{yx}^{(m)} p_{xy}^{(r)} \cdot \infty = \infty.
    \]
    Dunque anche lo stato $y$ è ricorrente.
\end{enumerate}
\end{dimostrazione}

\begin{oralbox}
``Incastriamo la serie di $p_{xx}$ dentro quella di $p_{yy}$ moltiplicando per i cammini fissi $y \to x$ e $x \to y$: se la prima diverge, costringe alla divergenza anche la seconda.''
\end{oralbox}

\begin{tipbox}
Formula da scrivere: \textbf{$p_{yy}^{(m+n+r)} \ge p_{yx}^{(m)} p_{xx}^{(n)} p_{xy}^{(r)}$}.
\end{tipbox}
\end{teoremabox}

\vspace{0.3cm}

\begin{teoremabox}{Teorema 7.8.9 --- Unicità della Misura Invariante}
Sia $X$ una catena di Markov irriducibile e ricorrente. Se esiste una misura invariante $\pi$ ($\pi P = \pi$), essa è unica a meno di un fattore moltiplicativo scalare. Se $\sum \pi(x) = 1$, la distribuzione stazionaria è strettamente unica.

\begin{cosafa}
Garantisce che una volta raggiunto l'equilibrio stazionario su una catena connessa e ricorrente, non possono esistere configurazioni di equilibrio alternative.
\end{cosafa}

\begin{dimostrazione}
Fissato uno stato di riferimento $z$, si costruisce la misura invariante $\mu_z(y) = \mathbb{E}_z\left[\sum_{n=0}^{\tau_z-1} \mathbf{1}_{\{X_n=y\}}\right]$ (tempo medio trascorso in $y$ prima di tornare a $z$). Qualsiasi altra misura invariante $\nu$ deve soddisfare per ricorsione $\nu(y) = \nu(z) \mu_z(y)$. Se normalizzata a somma $1$, il fattore $\nu(z)$ è univocamente determinato.
\end{dimostrazione}

\begin{oralbox}
``Ogni misura invariante è proporzionale al tempo medio di permanenza negli stati durante un ciclo di ritorno; normalizzando a 1 la costante si fissa univocamente.''
\end{oralbox}
\end{teoremabox}

\newpage

\begin{teoremabox}{Teorema 7.8.10 --- Equivalenza per Stati Ricorrenti Positivi}
Sia $X$ una catena di Markov irriducibile. Le seguenti affermazioni sono equivalenti:
\begin{enumerate}
    \item Esiste almeno uno stato $x$ ricorrente positivo ($\mathbb{E}_x[\tau_x] < \infty$).
    \item La catena ammette una misura di probabilità invariante $\pi$ ($\pi P = \pi$).
    \item Tutti gli stati $y \in E$ sono ricorrenti positivi.
\end{enumerate}

\begin{cosafa}
Collega la stazionarietà all'esistenza di tempi medi di ritorno finiti: avere equilibrio statistico globale equivale a dire che il sistema ritorna rapidamente in ogni suo punto.
\end{cosafa}

\begin{dimostrazione}
\begin{itemize}
    \item $(1 \implies 2)$: Se $\mu_x = \mathbb{E}_x[\tau_x] < \infty$, la misura di ciclo $\mu(y)$ ha massa totale $\sum_y \mu(y) = \mu_x < \infty$. Normalizzando $\pi(y) = \mu(y)/\mu_x$ otteniamo una probabilità invariante.
    \item $(2 \implies 3)$: Data una probabilità invariante $\pi$, per ogni stato $y$ vale la relazione fondamentale $\pi(y) = \frac{1}{\mathbb{E}_y[\tau_y]}$. Poiché $\pi(y) > 0$ per irriducibilità, il tempo atteso $\mathbb{E}_y[\tau_y] = 1/\pi(y) < \infty$, quindi ogni stato è ricorrente positivo.
\end{itemize}
\end{dimostrazione}

\begin{oralbox}
``Dimostriamo il ciclo: la misura di ciclo ha somma finita $\mathbb{E}_x[\tau_x]$, normalizzandola otteniamo $\pi$, e dal legame fondamentale $\pi(y) = 1/\mathbb{E}_y[\tau_y]$ deduciamo che tutti i tempi di ritorno sono finiti.''
\end{oralbox}

\begin{tipbox}
Formula magica da stampare in mente: \textbf{$\pi(y) = \frac{1}{\mathbb{E}_y[\tau_y]}$}.
\end{tipbox}
\end{teoremabox}

\vspace{0.3cm}

\begin{teoremabox}{Teorema 7.9.1 --- Teorema Ergodico per Catene di Markov}
Sia $X$ una catena di Markov irriducibile e ricorrente positiva con distribuzione stazionaria $\pi$. Per ogni funzione $f: E \to \mathbb{R}$ con $\sum |f(x)|\pi(x) < \infty$:
\[
\lim_{n \to \infty} \frac{1}{n} \sum_{k=1}^n f(X_k) = \sum_{x \in E} f(x) \pi(x) \quad \text{quasi certamente}.
\]

\begin{cosafa}
``La media temporale eguaglia la media spaziale'': calcolare la media dei valori raccolti lungo una singola traiettoria infinita dà lo stesso identico risultato che calcolare la media pesata con la distribuzione di equilibrio $\pi$.
\end{cosafa}

\begin{dimostrazione}
Spezziamo la traiettoria temporale in blocchi indipendenti tra i successivi ritorni a uno stato fissato $x_0$ (escursioni). Per la Proprietà Forte di Markov, la quantità $\sum f(X_k)$ calcolata all'interno di ciascuna escursione costituisce una successione di v.a. i.i.d. con media finita. Applicando la Legge Forte dei Grandi Numeri sui blocchi di escursione, il rapporto tra somma accumulata e tempo totale converge al valore atteso rispetto alla misura stazionaria.
\end{dimostrazione}

\begin{oralbox}
``Decomponiamo la traiettoria nei blocchi di ritorno a uno stato: per la proprietà forte le escursioni sono i.i.d., e applicando la Legge Forte sui blocchi ricaviamo l'uguaglianza tra media temporale e media spaziale $\mathbb{E}_\pi[f]$.''
\end{oralbox}

\begin{tipbox}
Parole d'oro: \textbf{Escursioni i.i.d. + Legge Forte dei Grandi Numeri sui blocchi}.
\end{tipbox}
\end{teoremabox}

\end{document}
